\chapter{KẾT LUẬN VÀ HƯỚNG PHÁT TRIỂN}

\section{Kết luận}

Qua quá trình nghiên cứu và thực hiện khóa luận, em đã hoàn thành được các mục tiêu đề ra ban đầu trên cả hai phương diện nghiên cứu lý thuyết và xây dựng sản phẩm phần mềm.

\subsection{Về nghiên cứu lý thuyết}

Khóa luận đã hệ thống hóa kiến thức chuyên sâu về kiến trúc mạng lưới Avalanche, đặc biệt là giao thức đồng thuận Snowball --- một hướng tiếp cận khác biệt hoàn toàn so với các giao thức BFT truyền thống. Thông qua phân tích toán học dựa trên mô hình chuỗi Markov, em đã chỉ ra cơ chế ``metastability'' giúp Snowball đạt được sự hội tụ nhanh với xác suất đảo ngược quyết định cực thấp ($< 10^{-9}$). Đồng thời, khóa luận cũng đã phân tích và so sánh mô hình Subnet của Avalanche với các giải pháp tương tự (Cosmos Zone, Polkadot Parachain), chỉ ra ưu thế về tính linh hoạt và chi phí thấp của Avalanche Subnet.

\subsection{Về sản phẩm phần mềm}

Hệ thống phần mềm được xây dựng đã giải quyết trực tiếp ba nhóm thách thức thực tế được xác định ở Chương 1:

\begin{enumerate}
    \item \textbf{Giảm độ phức tạp triển khai:} Web Wizard cho phép người dùng tạo và khởi động một Subnet hoàn chỉnh trong dưới 2 phút, giảm 95\% so với phương thức CLI thủ công (trung bình 42 phút). Toàn bộ quy trình phức tạp (tạo khóa, cấu hình Genesis, đăng ký validator, đồng bộ mạng) được tự động hóa hoàn toàn.

    \item \textbf{Cung cấp giám sát trực quan:} Dashboard real-time hiển thị đầy đủ các chỉ số vận hành quan trọng bao gồm block height, TPS, trạng thái node, peer count, CPU/RAM usage với độ trễ cập nhật dưới 3 giây thông qua WebSocket.

    \item \textbf{Giảm gánh nặng vận hành:} Tự động hóa các tác vụ DevOps bằng Docker container quản lý qua API: khởi động/dừng node, cấp phát port, thu thập metrics, xử lý sự cố. Hệ thống xác thực JWT + Google OAuth và phân quyền ownership đảm bảo tính bảo mật.
\end{enumerate}

\subsection{Đóng góp của đề tài}

Đề tài đóng góp chính vào ba lĩnh vực:

\begin{itemize}
    \item \textbf{Về kỹ thuật:} Đề xuất mô hình Orchestrator Pattern cho quản lý vòng đời container hóa validator node, kết hợp WebSocket cho giám sát real-time --- có thể áp dụng cho các hệ thống quản trị hạ tầng phân tán tương tự.

    \item \textbf{Về ứng dụng:} Sản phẩm phần mềm có tiềm năng phát triển thành giải pháp BaaS thương mại, phục vụ thị trường đang tăng trưởng mạnh (CAGR 62,2\% theo MarketsandMarkets \cite{marketsandmarkets2021baas}).

    \item \textbf{Về học thuật:} Khóa luận cung cấp tài liệu tham khảo hệ thống bằng tiếng Việt về kiến trúc Avalanche và giao thức Snowball, phục vụ cho các sinh viên và nhà nghiên cứu quan tâm đến lĩnh vực Blockchain Layer 1.
\end{itemize}

\section{Hướng phát triển}

Dựa trên những hạn chế đã nhận diện và xu hướng phát triển công nghệ, em đề xuất năm hướng mở rộng cho các nghiên cứu tiếp theo:

\subsection{Triển khai phân tán đa máy chủ (Multi-Host Deployment)}

Hiện tại hệ thống giới hạn trên một máy chủ duy nhất. Hướng phát triển quan trọng nhất là hỗ trợ triển khai validator node trên nhiều máy chủ vật lý hoặc cloud instance (AWS EC2, Google Cloud Compute Engine) thông qua Docker Swarm hoặc Kubernetes. Điều này cho phép mở rộng quy mô validator lên hàng trăm node và chịu tải tốt hơn.

\subsection{Auto-Scaling và Self-Healing}

Tích hợp cơ chế tự động mở rộng (auto-scaling) dựa trên metrics: khi tải mạng tăng cao (TPS tăng, latency tăng), hệ thống tự động thêm validator node; khi tải giảm, thu hồi tài nguyên. Kết hợp self-healing: khi phát hiện node bị crash (qua health check), tự động khởi động lại container hoặc thay thế bằng node mới.

\subsection{Tích hợp Avalanche Warp Messaging (AWM)}

Avalanche Warp Messaging là giao thức giao tiếp liên chuỗi (cross-chain communication) mới của Avalanche, cho phép các Subnet gửi thông điệp và chuyển tài sản qua lại mà không cần bridge bên thứ ba. Tích hợp AWM vào hệ thống sẽ mở ra khả năng xây dựng các ứng dụng đa chuỗi (multi-chain applications).

\subsection{Thương mại hóa dưới dạng SaaS}

Phát triển mô hình kinh doanh SaaS (Software-as-a-Service) với các gói dịch vụ phân tầng: Free (1 Subnet, giới hạn tài nguyên), Pro (5 Subnet, hỗ trợ ưu tiên), Enterprise (không giới hạn, SLA cam kết uptime 99,9\%). Tích hợp hệ thống thanh toán, quản lý quota và billing tự động.

\subsection{Ứng dụng Machine Learning cho giám sát thông minh}

Áp dụng các mô hình Machine Learning phân tích dữ liệu time-series từ Prometheus để dự đoán sự cố trước khi xảy ra (anomaly detection, predictive monitoring). Ví dụ: phát hiện xu hướng tăng RAM usage bất thường để cảnh báo trước khi node OOM (Out of Memory) crash.

\vspace{1cm}
\noindent\rule{\textwidth}{0.4pt}

\vspace{0.5cm}
Với những kết quả đã đạt được, em tin rằng đề tài không chỉ có giá trị trong phạm vi học thuật mà còn có tiềm năng ứng dụng thực tiễn cao. Hệ thống đã chứng minh tính khả thi của việc ``bình dân hóa'' công nghệ quản trị Blockchain --- biến một quy trình phức tạp đòi hỏi chuyên gia thành một trải nghiệm đơn giản, trực quan cho mọi người dùng.
